%% kapa: exam=1 round=37 year=2026 period=2 form=A subject=appraisal_law subject_ko=감정평가관계법규 seq=11 number=11
%% source: Q-Net 감정평가사 기출문제 (cst00309) artlSeq=5258924 file=raw/1008/5258924/2026년 제37회 감정평가사 1차 2교시.pdf
%% conversion: kordoc 4.24.51 -> tools/gen1.py (auto); review=proofread-sonnet; flags=none
\begin{question}{11}
국토의 계획 및 이용에 관한 법령상 시범도시의 지정ㆍ지원 등에 관한 설명으로 옳지 않은 것은?
\begin{choices}
  \item 시ㆍ도지사는 국토교통부장관에게 시범도시의 지정을 요청하고자 하는 때에는 미리 당해 시ㆍ도도시계획위원회의 자문을 거쳐야 한다.
  \item 국토교통부장관은 시범도시에 대하여 시범도시사업의 시행에 소요되는 비용(보상비 포함)의 50퍼센트 이하의 범위에서 보조 또는 융자를 할 수 있다.
  \item 시범도시의 공모에 응모할 수 있는 자는 특별시장ㆍ광역시장ㆍ특별자치시장ㆍ특별자치도지사ㆍ시장ㆍ군수 또는 구청장으로 한다.
  \item 관계 중앙행정기관의 장 또는 시ㆍ도지사는 국토교통부장관에게 시범도시의 지정을 요청하고자 하는 때에는 미리 설문조사ㆍ열람 등을 통하여 주민의 의견을 들은 후 관계 지방자치단체의 장의 의견을 들어야 한다.
  \item 국토교통부장관은 시범도시를 지정하려면 중앙도시계획위원회의 심의를 거쳐야 한다.
\end{choices}
\end{question}
\answer{2}
